\documentclass[11pt]{article}
\usepackage[margin=1in]{geometry}
\usepackage{amsmath,amssymb,mathtools,bm}
\usepackage{microtype}
\usepackage{lmodern}
\usepackage[T1]{fontenc}
\usepackage{hyperref}
\usepackage{enumitem}
\usepackage{booktabs}
\usepackage{tikz}
\usetikzlibrary{arrows.meta,positioning}
\setlength{\parindent}{0pt}
\setlength{\parskip}{0.65em}
\hypersetup{colorlinks=true,linkcolor=black,urlcolor=black,citecolor=black}
\newcommand{\FH}{\mathbf F_{\!H}}
\newcommand{\ket}[1]{\left|#1\right\rangle}
\newcommand{\bra}[1]{\left\langle#1\right|}
\newcommand{\braket}[2]{\left\langle#1\middle|#2\right\rangle}
\newenvironment{plainbox}[1]{\begin{center}\begin{minipage}{0.91\textwidth}\hrule\vspace{0.55em}\textbf{#1}\par\smallskip}{\vspace{0.55em}\hrule\end{minipage}\end{center}}

\newcommand{\AU}{\mathrm{AU}}
\newenvironment{model}[1][]{\par\medskip\noindent\textbf{Model#1.}\itshape\quad}{\upshape\par\medskip}

\newcommand{\Tr}{\operatorname{Tr}}
\newcommand{\tr}{\operatorname{tr}}
\title{\textbf{Hysterical Response from Quantum Alternatives}\\[0.45em]
\large Few-Particle Thought Experiments, a Solvable Two-Level Toy,\\Nonlocal Correlations, and the Conjectured Macroscopic Survival of Gravity\\[0.9em]
\normalsize Working conceptual formulation}
\author{Andrew Bruce Boyd}
\date{September 2026}

\begin{document}
\maketitle
\noindent\textit{Computational note.}
Exploratory experiments and paper writing were assisted by generative AI.
This primer states no formal mathematical theorems for Lean verification.

\begin{abstract}
Quantum mechanics allows physical systems to occupy superpositions of alternatives. In realistic macroscopic systems, interaction with the environment rapidly decoheres those alternatives, making them behave as effectively separate outcomes. Decoherence is extraordinarily effective, but ``effectively zero'' is not automatically the same statement as ``exactly zero.'' This motivates a simple question: could an open quantum system retain an extremely small dynamical response to alternative interaction histories that have been almost, but not perfectly, erased from its local history?

We call the hypothetical residual, history-dependent effect the \textbf{Hysterical Response}. This paper is intentionally nontechnical. It does not derive the effect from a fundamental theory. Instead, two-, three-, and four-particle thought experiments, together with a solvable two-level gravitational toy, are used to explain what such a response would mean, why correlations make it intrinsically nonlocal at the level of the joint quantum state, and why a macroscopic response need not be a simple sum of independent particle contributions. We conjecture that every fundamental interaction can exhibit microscopic Hysterical Response, while only the gravitational response remains macroscopically appreciable in an appropriate many-particle limit. These conjectures are motivations, not results. The purpose of this primer is to make the physical idea clear before the technical papers ask whether it survives mathematics and observation.
\end{abstract}
\tableofcontents
\newpage

\section{The question}
Consider two particles, $A$ and $B$. In ordinary Newtonian gravity, if their masses are $m_A$ and $m_B$ and their separation is $r$, the attraction has magnitude
\begin{equation}
F_N=\frac{Gm_Am_B}{r^2}.
\end{equation}
Nothing unusual has happened yet.

Now suppose $B$ is placed in a spatial superposition,
\begin{equation}
\ket{B_L}+\ket{B_R}.
\end{equation}
This is the same basic idea expressed by the familiar shorthand
\begin{equation}
\ket{\text{cat alive}}+\ket{\text{cat dead}}.
\end{equation}
Real objects do not normally appear as visible superpositions because they interact continuously with their surroundings. Photons scatter, air molecules collide, heat is exchanged, and information about the alternatives spreads into the environment. Schematically,
\begin{equation}
\ket{B_L}\ket{E_L}+\ket{B_R}\ket{E_R}.
\end{equation}
For a strongly decohered system,
\begin{equation}
\braket{E_R}{E_L}\approx0.
\end{equation}
The symbol $\approx$ is the point of departure. What if an open system can retain an extraordinarily small dynamical memory of alternatives that have become effectively, but not perfectly, separated?

\section{The Hysterical Response}
\begin{plainbox}{Conjectural phenomenon: Hysterical Response}
An open quantum system may retain an extremely small dynamical response to alternative interaction histories that have been strongly, but not perfectly, decohered.
\end{plainbox}

The name is intentional. Ordinary \emph{hysteresis} describes systems whose present behaviour depends partly on their past. Here we use the deliberately unconventional adjective \emph{Hysterical} for a proposed response whose present value retains information about quantum alternatives in the system's interaction history.

We denote an effective Hysterical force contribution by $\FH$ and write phenomenologically
\begin{equation}
\boxed{\mathbf F_{\rm observed}=\mathbf F_{\rm ordinary}+\FH.}
\end{equation}
This is not proposed as a fundamental quantum force operator. It is the quantity that an observer using ordinary dynamics would interpret as an additional force. The technical companion asks whether such a response can arise consistently from open-system dynamics; the present primer only explains the idea.

\section{Toy model I: two particles}
Let $A$ be a probe and $B$ a source. Suppose $B$ has a near alternative $B_N$ and a far alternative $B_F$, at distances $r_N<r_F$. Then
\begin{equation}
F_N=\frac{Gm_Am_B}{r_N^2},\qquad F_F=\frac{Gm_Am_B}{r_F^2},\qquad F_N>F_F.
\end{equation}
Suppose the experienced configuration contains $B_N$. Ordinary reasoning gives $F_N$. The Hysterical conjecture asks whether the alternative interaction history involving $F_F$ can leave a tiny residual response,
\begin{equation}
F_{\rm observed}=F_N+F_H.
\end{equation}
As a mnemonic only, one might imagine
\begin{equation}
F_H\sim\varepsilon\,(F_N-F_F),\qquad |\varepsilon|\ll1,
\end{equation}
where $\varepsilon$ represents an extremely small surviving memory. This is deliberately \emph{not} a microscopic law.

The two-particle example gives the first intuition:
\begin{equation}
\boxed{\text{different quantum interaction histories}\ \Longrightarrow\ \text{possible residual dynamical response}.}
\end{equation}
But it also makes the effect look more local and additive than it need be.

\section{Toy model II: three particles and the correlation effect}
Now let $A$ remain the probe while $B$ and $C$ are sources. Each can be near or far from $A$. Compare
\begin{align}
\ket{\Psi_+}&=\frac{1}{\sqrt2}\left(\ket{B_NC_N}+\ket{B_FC_F}\right),\\
\ket{\Psi_-}&=\frac{1}{\sqrt2}\left(\ket{B_NC_F}+\ket{B_FC_N}\right).
\end{align}
If we look at $B$ alone, both states give
\begin{equation}
P(B_N)=P(B_F)=\frac12.
\end{equation}
The same is true for $C$:
\begin{equation}
P(C_N)=P(C_F)=\frac12.
\end{equation}
So the individual behaviour of $B$ is the same in both experiments, and the individual behaviour of $C$ is the same. What changes is their \emph{correlation}.

Assume equal source masses and a symmetric geometry. Let one near source contribute force $F_N$ and one far source $F_F$. In $\ket{\Psi_+}$ the alternatives have total forces
\begin{equation}
F_{NN}=2F_N,\qquad F_{FF}=2F_F,
\end{equation}
so their contrast is
\begin{equation}
\Delta F_+=2(F_N-F_F).
\end{equation}
For $\ket{\Psi_-}$,
\begin{equation}
F_{NF}=F_N+F_F,\qquad F_{FN}=F_F+F_N,
\end{equation}
and therefore
\begin{equation}
\boxed{\Delta F_-=0.}
\end{equation}

This is the central motivating example. Two states with identical one-particle statistics can expose $A$ to very different \emph{joint} gravitational histories. A Hysterical Response that remembers those histories can therefore satisfy
\begin{equation}
\boxed{F_H^{(A)}[\Psi_+]\neq F_H^{(A)}[\Psi_-],}
\end{equation}
even though the separate statistics of $B$ and $C$ have not changed.

\begin{plainbox}{The ``oh, I get it'' point}
The Hysterical Response cannot in general be just ``a tiny extra force from every unrealized particle position.'' If it remembers alternative interaction histories, it must be capable of remembering how the alternatives of different particles are correlated.
\end{plainbox}

\subsection{Operator form and basis independence}
Write $|N\rangle\equiv|0\rangle$ and $|F\rangle\equiv|1\rangle$, and introduce the near/far contrast operator
\begin{equation}
Z=|N\rangle\langle N|-|F\rangle\langle F|.
\end{equation}
Then both states have vanishing one-body contrasts $\langle Z_B\rangle=\langle Z_C\rangle=0$, while
\begin{equation}
\boxed{\langle Z_BZ_C\rangle_{+}=+1,\qquad \langle Z_BZ_C\rangle_{-}=-1.}
\end{equation}
A correlation-sensitive probe observable of the schematic form
\begin{equation}
R_A=r_0I+r_BZ_B+r_CZ_C+r_{BC}Z_BZ_C+\cdots
\end{equation}
therefore obeys
\begin{equation}
\boxed{\langle R_A\rangle_{+}-\langle R_A\rangle_{-}=2r_{BC}}
\end{equation}
when those terms are retained.

The same physics looks different in another basis. With
\begin{equation}
|+\rangle=\frac{|N\rangle+|F\rangle}{\sqrt2},\qquad
|-\rangle=\frac{|N\rangle-|F\rangle}{\sqrt2},
\end{equation}
one finds
\begin{align}
\ket{\Psi_+}&=\frac{|++\rangle+|--\rangle}{\sqrt2},\\
\ket{\Psi_-}&=\frac{|++\rangle-|--\rangle}{\sqrt2},
\end{align}
so the distinction appears as a relative phase among Hadamard configurations, while $Z\mapsto X$ and
\begin{equation}
\langle X_BX_C\rangle_{+}=+1,\qquad \langle X_BX_C\rangle_{-}=-1.
\end{equation}
The response difference $2r_{BC}$ is unchanged: a passive change of representation cannot create or destroy the physical discriminator.
This is the concrete reason later technical work insists that ``off-diagonal'' language is always relative to a physically selected classicalization map, not to an arbitrary Hilbert-space coordinate basis.

\section{What we mean by nonlocal}
The three-particle example immediately makes the response nonlocal in an important but limited sense. The response at $A$ need not be determined by independent local descriptions of $B$ and $C$; it may depend on their joint state:
\begin{equation}
\boxed{\FH(A)=\mathcal F_H[\text{joint state},\ \text{correlations},\ \text{environmental history}].}
\end{equation}
Changing correlations among spatially separated components can therefore change the Hysterical response at $A$ even when their individual reduced descriptions remain unchanged.

This is \emph{state nonlocality}. We do not claim here that it permits faster-than-light signalling, nor that relativistic causality is violated. Those are stronger questions requiring a relativistic theory that has not yet been constructed.

\section{Toy model III: four particles and higher correlations}
Three particles show that pair correlations can matter. A fourth raises the possibility that still higher joint structure matters.

Let $A$ be the probe and $B,C,D$ the source system. The response might admit a schematic hierarchy
\begin{equation}
F_H=F_H^{(1)}+F_H^{(2)}+F_H^{(3)}+\cdots,
\end{equation}
where the superscripts denote information carried by one-source structure, pair correlations, irreducible three-source correlations, and so on. We do not assert that this expansion is fundamental; it simply illustrates why the many-particle response need not reduce to a sum over independent pairs.

The lesson is that the microscopic Hysterical Response may be \textbf{configuration-sensitive}, rather than merely \textbf{density-sensitive}.

\section{Why should the effect be gravitational?}
Nothing in the motivating argument so far knows that we chose gravity. Replace the gravitational interaction by electrostatics,
\begin{equation}
V_{\rm EM}(r)=\frac{q_Aq_B}{4\pi\epsilon_0r}.
\end{equation}
A spatial superposition again corresponds to alternative interaction energies and forces. If Hysterical Response originates from quantum alternatives and environmental memory, rather than from a special property of gravity, it is natural to ask whether electromagnetism should exhibit it too.

\begin{plainbox}{Conjecture 1: Interaction universality}
Every fundamental interaction can, in principle, exhibit a microscopic Hysterical Response when an open quantum system contains incompletely erased alternative interaction histories.
\end{plainbox}

Symbolically,
\begin{equation}
\boxed{\mathbf F_X^{\rm observed}=\mathbf F_X^{\rm ordinary}+\mathbf F_{H,X},}
\end{equation}
where $X$ labels an interaction channel. This is a conjecture. A microscopic theory could instead show that gravity or stress-energy is special.

\section{Then why do we not see Hysterical electromagnetism everywhere?}
Electromagnetism is enormously stronger than gravity between elementary charged particles. If its Hysterical contribution survived straightforwardly into bulk matter, that would be an immediate problem.

But ordinary matter is approximately electrically neutral. Positive and negative charge can cancel:
\begin{equation}
(+q)+(-q)=0.
\end{equation}
Ordinary Newtonian mass does not have the analogous cancellation in normal matter: positive mass contributions add.

This gives an intuitive reason why a microscopic effect shared by several interactions might nevertheless appear gravitational on large scales. It does not prove it.

Suppose a positive and negative charge acquire Hysterical weights $h_+$ and $h_-$. Their residual contribution could contain
\begin{equation}
q h_+-q h_-.
\end{equation}
Charge neutrality says $q-q=0$; it does not imply $qh_+-qh_-=0$ unless the response dynamics enforce $h_+=h_-$. Because the weights themselves may depend on nonlocal correlations and history, neutrality alone is insufficient.

\section{Toy model IV: a solvable two-level gravitational block}
The particle examples say \emph{what} a residual response would mean. A deliberately minimal channel shows \emph{how} a leftover force can sit on an ordinary interaction after decoherence, without inventing a new force carrier.

Imagine just two localized alternatives---``near'' versus ``far,'' or two pointer placements the environment treats as classical.
Write a two-level Hamiltonian and an ordinary force that differs slightly between them,
\begin{equation}
H=-\kappa X+\frac{\Delta}{2}Z,\qquad
F=\bar f\,I+\frac{\delta f}{2}Z,\qquad
\Omega=\sqrt{4\kappa^2+\Delta^2}.
\end{equation}
Here $Z$ labels the two alternatives, so $\delta f$ is their \emph{force contrast}.
Thermal relaxation toward the equilibrium population of $Z$ produces a \emph{population} force bias (not an off-diagonal ``coherence force'' relative to localization),
\begin{equation}
B=-\frac{\delta f\,\Delta}{2\Omega}\tanh\Bigl(\frac{\beta\Omega}{2}\Bigr).
\end{equation}
Randomly phased arrivals wipe mean coherence in the ensemble; what can survive is this population effect.

Refresh the channel by Poisson arrivals at rate $J$, let friction align populations at rate $\Gamma_{\mathrm{fr}}$, and let emission or removal cut the lifetime at rate $\Gamma_{\mathrm{em}}$.
Under an independent-renewal picture the steady leftover force factors as
\begin{equation}
F^{\mathrm{ss}}_{\mathrm{diss}}
=\frac{J}{\Gamma_{\mathrm{em}}}
\frac{\Gamma_{\mathrm{fr}}}{\Gamma_{\mathrm{em}}+\Gamma_{\mathrm{fr}}}\,B.
\end{equation}
So the leftover is a renewal prefactor times a bias built from force contrast and energetic splitting.
The same elementary algebra yields an inheritance bound:
\begin{equation}
\boxed{|B|\le\frac{|\delta f|}{2}.}
\end{equation}
Killing the force contrast between the two alternatives kills the bias that sources the leftover---with no further restriction on splitting or temperature.

\begin{plainbox}{What the two-level toy is for}
It is a thought-experiment calculator, not a claim that Nature is two-dimensional.
It separates population bias from coherence, shows how arrivals/friction/emission keep a channel alive, and makes precise why neutrality of the \emph{composite} force contrast matters more than summing independent little biases.
\end{plainbox}

This connects to the earlier configuration-versus-density lesson.
The bias $B$ cares about the contrast between alternatives---a configuration difference---not merely about a mean mass density.
Two arrangements with the same average density but different force contrasts can therefore leave different leftovers, just as $\ket{\Psi_+}$ and $\ket{\Psi_-}$ did with identical one-particle statistics.
Opposite-sign constituents illustrate the same point about order: forming $\delta f_{\mathrm{comp}}=\delta f_++\delta f_-$ first can cancel; adding independent biases $B(q)+B(-q)$ need not.

\section{The gravitational survival conjecture}
We therefore separate two claims: microscopic existence and macroscopic survival.

\begin{plainbox}{Conjecture 2: Gravitational survival}
Under the appropriate macroscopic coarse-graining and many-particle limit, nongravitational Hysterical responses vanish or become macroscopically negligible, while the gravitational Hysterical Response remains macroscopically appreciable.
\end{plainbox}

Very schematically,
\begin{equation}
N\to\infty:\qquad F_{H,g}\ \text{survives},\qquad F_{H,X}\to0\ \text{macroscopically for }X\neq g.
\end{equation}
We do not yet define the correct thermodynamic or coarse-grained limit, and we do not assume a particular $N$-scaling. For electromagnetism, signed charge cancellation is suggestive; for the strong interaction, confinement is suggestive; for the weak interaction, short range is suggestive. None is yet a derivation of the Hysterical many-body limit.

\section{From a few particles to a galaxy}
For one microscopic system, any residual response after decoherence could be fantastically small. An astronomical open system is different in scale and structure. A galaxy contains enormous numbers of particles, continually forms new structures, exchanges matter with its surroundings, emits radiation, and continually creates new environmental records.

The motivating possibility is
\begin{equation}
\boxed{\text{tiny residual response per alternative}\times\text{enormous effective multiplicity}\longrightarrow\text{finite macroscopic response}.}
\end{equation}
This is qualitative. It is not a literal count of Everett branches, and it does not assume that microscopic contributions add independently. The three-particle model already tells us that correlations can change the answer.

A macroscopic response variable must therefore be allowed to summarize more than local mass density. Schematically one might eventually require
\begin{equation}
Q=Q[\text{matter},\ \text{correlations},\ \text{environment},\ \text{history}].
\end{equation}
Later papers use an effective response variable of this kind without claiming that its microscopic form has already been derived.

\section{Where does the system end?}
Nonlocal correlation sensitivity creates a boundary problem. If we draw a sphere of radius $R$ around a galaxy and call everything inside ``the galaxy,'' why should the Hysterical Response care about our bookkeeping choice? Matter outside the chosen region may remain relevant through correlations or through the response history.

A physically meaningful quantity might eventually involve a large-region limit such as
\begin{equation}
\lim_{R\to\infty}Q(R),
\end{equation}
if that limit exists, or a dynamically generated memory kernel that suppresses sufficiently distant contributions. This primer does not choose between these possibilities. The point is simply that once correlations matter, defining the relevant system becomes a physical question.

\section{What has been shown, assumed, and conjectured}
To keep the logic explicit, the status of the claims in this paper is as follows.

\textbf{Elementary mathematical observation.} Two joint states can have identical one-particle marginals while possessing different correlations. In the symmetric three-particle example above, those correlations produce different branch-by-branch gravitational force contrasts at $A$. In the two-level toy, the leftover-sourcing population bias obeys $|B|\le|\delta f|/2$.

\textbf{Phenomenological postulate.} A residual history-dependent force-like quantity, called the Hysterical Response, is allowed to depend on alternative interaction histories.

\textbf{Conjecture.} The Hysterical Response can depend on joint correlations and environmental history, making it nonlocal in the joint-state sense described above.

\textbf{Conjecture.} Every fundamental interaction can exhibit microscopic Hysterical Response.

\textbf{Conjecture.} Only gravity remains macroscopically appreciable in the appropriate $N\to\infty$ limit.

We have not derived the Hysterical Response from quantum mechanics; proved that residual environmental overlap generates an observable force; proved electromagnetic cancellation; established the many-particle limit; constructed a relativistic completion; or shown that the effect explains dark-matter phenomenology.

\section{Why call it Hysterical?}
The proposed phenomenon is fundamentally about memory. Its present value is conjectured not to depend solely on the instantaneous classical configuration, but also on alternative configurations through which the quantum state and its environment have evolved.

In ordinary hysteresis, history matters. Here, quantum history matters. Hence
\begin{equation}
\boxed{\textbf{Hysterical Response}.}
\end{equation}
A corresponding macroscopic field will be called the \textbf{Hysterical Response field}, or simply the \textbf{Hysterical field} where the meaning is clear.

\section{The conceptual picture}
The proposal can be summarized without technical machinery:
\begin{equation}
\begin{gathered}
\text{quantum alternatives}\\
\downarrow\\
\text{environmental decoherence}\\
\downarrow\\
\text{incomplete dynamical erasure, conjecturally}\\
\downarrow\\
\textbf{Hysterical Response}\\
\downarrow\\
\text{correlation-sensitive, nonlocal joint-state behaviour}\\
\downarrow\\
\text{possible macroscopic gravitational survival.}
\end{gathered}
\end{equation}

The two-particle example says what a residual response would mean. The three-particle example explains why correlations matter. The four-particle example warns that higher-order structure may matter too. The two-level gravitational block shows how a population bias controlled by force contrast can yield a steady leftover under arrivals, friction, and emission. Repeating the thought experiment with Coulomb interactions motivates microscopic interaction universality. The contrasting behaviour of signed charge and positive Newtonian mass motivates, but does not prove, gravitational macroscopic survival.

\section{What comes next}
The simplicity of this primer conceals a serious problem. One cannot simply declare that an off-diagonal part of a Newtonian force operator is ``force leaking between branches.'' Under natural localization assumptions that construction can vanish identically.

The technical companion therefore asks the first mathematical question: can a Hysterical Response arise through legitimate open-system dynamics, with coherent mixing and dissipation converting residual alternative structure into an ordinary population/response effect? The two-level toy above is the conceptual bridge to that question; the technical development treats the abstract resolvent response and neutrality theorems without needing the pedagogical walkthrough.

A subsequent many-body paper asks a different question: if microscopic Hysterical Response is interaction-general, which interaction channels actually survive coarse-graining as $N$ becomes large? In particular, can electromagnetic Hysterical Response disappear in neutral macroscopic matter without simply assuming the desired answer?

Only after those questions are separated does it make sense to investigate a surviving gravitational Hysterical field on galactic scales.

\section{Conclusion}
This primer proposes a simple physical possibility. Decoherence makes quantum alternatives effectively separate, but an open system may not erase their dynamical influence with infinite perfection. We call any resulting residual, history-dependent effect the \textbf{Hysterical Response}.

The most important motivating example is not the two-particle superposition but the three-particle correlation experiment. Two source systems can have identical individual particle statistics while possessing different joint alternatives and therefore different interaction histories at a probe. If the Hysterical Response remembers those histories, it must be capable of depending on correlations rather than only on independent local densities.

This motivates a response that is nonlocal in the joint-state sense. We conjecture that the microscopic mechanism is interaction-general, but separately conjecture that only gravity remains macroscopically appreciable in an appropriate many-particle limit. Neither statement is assumed to be proved by the examples.

The remainder of the series asks whether this intuitive picture survives increasingly demanding mathematics and observation.

\end{document}
